\documentclass[12pt,a4paper]{article}
\usepackage[utf8]{inputenc}
\usepackage[vietnamese]{babel}
\usepackage{amsmath,amssymb,amsthm}
\usepackage{geometry}
\usepackage{enumitem}
\usepackage[most]{tcolorbox}
\usepackage{hyperref}
\usepackage{fancyhdr}
\usepackage{tikz}
\usepackage{eso-pic}
\usepackage{xcolor}
\geometry{a4paper, margin=2cm, bottom=2.5cm}
\hypersetup{colorlinks=true, linkcolor=blue!70!black}
\setlist[itemize]{topsep=4pt, itemsep=2pt}
\setlist[enumerate]{topsep=4pt, itemsep=2pt}
\AddToShipoutPictureFG{%
  \AtPageCenter{%
    \begin{tikzpicture}[remember picture, overlay]
      \node[rotate=45, scale=1, opacity=0.06]{\fontsize{60}{72}\selectfont\bfseries\sffamily Đặng Tấn Quí};
    \end{tikzpicture}%
  }%
}
\pagestyle{fancy}
\fancyhf{}
\fancyhead[L]{\small\textit{LÝ THUYẾT SỐ}}
\fancyhead[R]{\small\textit{Đặng Tấn Quí}}
\fancyfoot[C]{\thepage}
\fancyfoot[L]{\footnotesize\textcopyright\ Đặng Tấn Quí}
\fancyfoot[R]{\footnotesize SĐT: 0377\,122\,210}
\renewcommand{\headrulewidth}{0.4pt}
\renewcommand{\footrulewidth}{0.4pt}
\fancypagestyle{plain}{\fancyhf{}\fancyfoot[C]{\thepage}\fancyfoot[L]{\footnotesize\textcopyright\ Đặng Tấn Quí}\fancyfoot[R]{\footnotesize SĐT: 0377\,122\,210}\renewcommand{\headrulewidth}{0pt}\renewcommand{\footrulewidth}{0.4pt}}
\newtcolorbox{dinhnghia}[1][]{enhanced, breakable, colback=blue!3!white, colframe=blue!60!black, fonttitle=\bfseries, title={Định nghĩa}, #1}
\newtcolorbox{dinhly}[1][]{enhanced, breakable, colback=green!3!white, colframe=green!50!black, fonttitle=\bfseries, title={Định lý}, #1}
\newtcolorbox{tinhchat}[1][]{enhanced, breakable, colback=yellow!5!white, colframe=orange!50!black, fonttitle=\bfseries, title={Tính chất}, #1}
\newtcolorbox{phuongphap}[1][]{enhanced, breakable, colback=gray!5!white, colframe=gray!50!black, fonttitle=\bfseries, title={Phương pháp}, #1}
\newtcolorbox{luuy}[1][]{enhanced, breakable, colback=red!3!white, colframe=red!50!black, fonttitle=\bfseries, title={Lưu ý}, #1}
\newtcolorbox{congthuc}[1][]{enhanced, breakable, colback=cyan!3!white, colframe=cyan!50!black, fonttitle=\bfseries, title={Công thức}, #1}
\newtcolorbox{vidu}[1][]{enhanced, breakable, colback=violet!3!white, colframe=violet!50!black, fonttitle=\bfseries, title={Ví dụ}, #1}

\begin{document}
\thispagestyle{empty}
\begin{tikzpicture}[remember picture, overlay]
  \fill[blue!80!black] (current page.south west) rectangle (current page.north east);
  \node[anchor=west, white, font=\fontsize{26}{32}\bfseries\selectfont]
    at ([xshift=3cm, yshift=-7.5cm]current page.north west) {LÝ THUYẾT SỐ};
  \node[anchor=west, white, font=\fontsize{18}{24}\selectfont]
    at ([xshift=3cm, yshift=-9cm]current page.north west) {Chia hết, đồng dư, Euler--Fermat, thặng dư bậc hai, hàm số học};
  \node[anchor=west, white, font=\Large\bfseries] at ([xshift=3cm, yshift=-13cm]current page.north west) {Biên soạn:};
  \node[anchor=west, yellow!80!white, font=\fontsize{22}{28}\bfseries\selectfont] at ([xshift=3cm, yshift=-14.5cm]current page.north west) {ĐẶNG TẤN QUÍ};
  \node[anchor=south east, white, font=\Large\bfseries] at ([xshift=-2cm, yshift=3cm]current page.south east) {\the\year};
\end{tikzpicture}
\newpage
\tableofcontents
\newpage
%% ================================================================
%%  CHƯƠNG 1: CHIA HẾT VÀ SỐ NGUYÊN TỐ
%% ================================================================
\section{Chia hết và số nguyên tố}

\begin{dinhnghia}[title={Chia hết}]
  $a \mid b$ ($a$ chia hết $b$): $\exists k \in \mathbb{Z}$: $b = ka$.
\end{dinhnghia}

\begin{dinhly}[title={Chia Euclid}]
  $\forall a, b \in \mathbb{Z}$, $b > 0$: $\exists!$ $q, r$: $a = bq + r$, $0 \le r < b$.
\end{dinhly}

\begin{dinhnghia}[title={ƯCLN, BCNN}]
  $\gcd(a, b)$: ước chung lớn nhất. $\text{lcm}(a, b)$: bội chung nhỏ nhất.

  $\gcd(a,b) \cdot \text{lcm}(a,b) = |ab|$.
\end{dinhnghia}

\begin{phuongphap}[title={Thuật toán Euclid}]
  $\gcd(a, b) = \gcd(b, a \bmod b)$. Lặp đến $r = 0$.

  \textbf{Euclid mở rộng}: tìm $x, y$: $ax + by = \gcd(a, b)$ (Bézout).
\end{phuongphap}

\begin{dinhly}[title={Định lý cơ bản của số học}]
  Mọi $n \ge 2$ phân tích duy nhất: $n = p_1^{e_1} p_2^{e_2} \cdots p_k^{e_k}$, $p_i$ nguyên tố.
\end{dinhly}

\begin{dinhly}
  Số lượng số nguyên tố là vô hạn (Euclid). \textbf{Định lý số nguyên tố:} $\pi(x) \sim \dfrac{x}{\ln x}$.
\end{dinhly}

%% ================================================================
%%  CHƯƠNG 2: ĐỒNG DƯ
%% ================================================================
\section{Đồng dư}

\begin{dinhnghia}
  $a \equiv b \pmod{m}$ nếu $m \mid (a - b)$. Quan hệ tương đương; $\mathbb{Z}_m = \mathbb{Z}/m\mathbb{Z}$ là vành.
\end{dinhnghia}

\begin{tinhchat}
  \begin{itemize}
    \item $a \equiv b$, $c \equiv d$ $\pmod{m}$ $\Rightarrow$ $a \pm c \equiv b \pm d$, $ac \equiv bd$.
    \item $a$ có nghịch đảo modulo $m$ $\Leftrightarrow$ $\gcd(a, m) = 1$.
  \end{itemize}
\end{tinhchat}

\begin{dinhly}[title={Định lý thặng dư Trung Hoa (CRT)}]
  Nếu $m_1, \ldots, m_k$ đôi một nguyên tố cùng nhau thì hệ:
  \[
    x \equiv a_i \pmod{m_i}, \quad i = 1, \ldots, k
  \]
  có nghiệm duy nhất modulo $M = m_1 m_2 \cdots m_k$.
\end{dinhly}

\begin{dinhnghia}[title={Hàm Euler}]
  $\varphi(n) = |\{k : 1 \le k \le n,\, \gcd(k, n) = 1\}|$.

  $\varphi(p^e) = p^{e-1}(p-1)$. Tính nhân: $\gcd(m,n)=1 \Rightarrow \varphi(mn) = \varphi(m)\varphi(n)$.
\end{dinhnghia}

\begin{dinhly}[title={Euler--Fermat}]
  Nếu $\gcd(a, m) = 1$: $a^{\varphi(m)} \equiv 1 \pmod{m}$.

  \textbf{Fermat nhỏ} ($m = p$ nguyên tố): $a^{p-1} \equiv 1 \pmod{p}$, hay $a^p \equiv a \pmod{p}$.
\end{dinhly}

\begin{dinhly}[title={Wilson}]
  $p$ nguyên tố $\Leftrightarrow$ $(p-1)! \equiv -1 \pmod{p}$.
\end{dinhly}

\begin{dinhnghia}[title={Bậc và căn nguyên thủy}]
  \textbf{Bậc} $\text{ord}_m(a)$: $d$ nhỏ nhất sao cho $a^d \equiv 1 \pmod{m}$. $\text{ord}_m(a) \mid \varphi(m)$.

  $g$ là \textbf{căn nguyên thủy} mod $m$ nếu $\text{ord}_m(g) = \varphi(m)$. Tồn tại khi $m = 2, 4, p^k, 2p^k$.
\end{dinhnghia}

%% ================================================================
%%  CHƯƠNG 3: ĐỒNG DƯ BẬC HAI VÀ THẶNG DƯ
%% ================================================================
\section{Đồng dư bậc hai}

\begin{dinhnghia}[title={Thặng dư bậc hai}]
  $a$ là \textbf{thặng dư bậc hai} mod $p$ nếu $\exists x$: $x^2 \equiv a \pmod{p}$. \textbf{Ký hiệu Legendre:}
  \[
    \left(\frac{a}{p}\right) = \begin{cases} 1 & \text{nếu } a \text{ là thặng dư bậc hai} \\ -1 & \text{nếu } a \text{ là phi thặng dư bậc hai} \\ 0 & \text{nếu } p \mid a. \end{cases}
  \]
\end{dinhnghia}

\begin{dinhly}[title={Tiêu chuẩn Euler}]
  $\left(\dfrac{a}{p}\right) \equiv a^{(p-1)/2} \pmod{p}$.
\end{dinhly}

\begin{dinhly}[title={Luật thuận nghịch bậc hai (Gauss)}]
  Với $p, q$ nguyên tố lẻ phân biệt:
  \[
    \left(\frac{p}{q}\right)\left(\frac{q}{p}\right) = (-1)^{\frac{p-1}{2}\cdot\frac{q-1}{2}}.
  \]
  Bổ sung: $\left(\dfrac{-1}{p}\right) = (-1)^{(p-1)/2}$, \;$\left(\dfrac{2}{p}\right) = (-1)^{(p^2-1)/8}$.
\end{dinhly}

%% ================================================================
%%  CHƯƠNG 4: HÀM SỐ HỌC
%% ================================================================
\section{Hàm số học}

\begin{dinhnghia}[title={Một số hàm số học quan trọng}]
  \begin{itemize}
    \item $\tau(n) = \sum_{d|n} 1$ (số ước).
    \item $\sigma(n) = \sum_{d|n} d$ (tổng ước).
    \item $\mu(n)$: hàm Möbius --- $\mu(1)=1$; $\mu(n) = (-1)^k$ nếu $n = p_1\cdots p_k$; $\mu(n) = 0$ nếu $n$ có ước bình phương.
  \end{itemize}
\end{dinhnghia}

\begin{dinhly}[title={Công thức nghịch đảo Möbius}]
  Nếu $f(n) = \sum_{d|n} g(d)$ thì $g(n) = \sum_{d|n} \mu(n/d)\,f(d)$.
\end{dinhly}

\begin{dinhly}[title={Tích chập Dirichlet}]
  $(f * g)(n) = \sum_{d|n} f(d)\,g(n/d)$. Vành giao hoán. $\varphi = \mu * \text{id}$, $\sigma = \text{id} * \mathbf{1}$.
\end{dinhly}

\end{document}